\chapter{RC 电路与 RL 电路}\label{ch:g12:rc-rl-circuits}

按下相机快门至一半：一声细微的嗡鸣持续两秒，
一盏灯变绿——只有此时闪光灯才会点亮，在一毫秒内释放
它用两秒钟积蓄的能量。本章将介绍使电路具有“记忆”与“时钟”功能的两种元件
——电容器与线圈——并
引入物理学中第一批微分方程，可用本年度所学的导数加以检验。

\section{电容器}

\begin{definition}[电容器与电容]\label{def:g12:rc-rl-circuits:capacitor}
\emph{电容器}\index{电容器}是一对相对放置的金属极板，
中间隔以薄的绝缘体。接入电源后，两极板分别带
等量异号的电荷 $+q$ 与 $-q$——两板之间的间隙中充满电场
（见\cref{ch:g11:electric-gravitational-fields}），两极板间的电压 $u$
与所储存的电荷成正比：
\[ q = C\,u . \]
常数 $C$ 称为\emph{电容}\index{电容}，其单位为
\emph{法拉}\index{法拉}（\unit{F}）：$\qty{1}{F} = \qty{1}{C/V}$。
\end{definition}

\begin{example}[数量级]\label{ex:g12:rc-rl-circuits:orders}
一法拉非常大，相当于每伏特储存一库仑电荷。无线电电路中的陶瓷电容器
：\unit{pF} 至 \unit{nF}；电源与
闪光灯中的电解电容器：\unit{\micro F} 至 \unit{mF}；超级电容器：可达数千法拉。
\end{example}

\begin{proposition}[流入电容器的电流]\label{prop:g12:rc-rl-circuits:cdudt}
流入电容器某一极板的电流，等于其电荷
增长的变化率：
\[
i = \frac{dq}{dt} = C\,\frac{du}{dt} ,
\]
小写字母表示随时间变化的量（$U$、$I$ 保持不变，
见\cref{ch:g11:circuits-and-power}）。
\end{proposition}

\begin{proof}
电流是单位时间内输送的电荷
（见\cref{def:g11:circuits-and-power:current}）：在极短的时间间隔内，
$i$ 即为 $q$ 对时间的导数；由 $q = Cu$，且 $C$ 为常数，
得 $C\,du/dt$。
\end{proof}

\begin{remark}[两个推论]\label{rem:g12:rc-rl-circuits:continuity}
若 $u$ 为常数，则 $i = 0$：已充电的电容器不会有稳定的电流通过。
此外，$u$ 不可能突变——若发生突变，则需要无穷大的电流：
电容器两端的电压即是电路的“记忆”。
\end{remark}

\section{电容器的充电：RC 电路}

一个回路由电源 $\mathcal{E}$（见\cref{def:g11:circuits-and-power:emf}）、
开关、电阻器与不带电的电容器组成：闭合 $K$ 的时刻为 $t = 0$。

\begin{omfigure}
\begin{circuitikz}[font=\small]
  \draw (0,0) to[battery1, l=$\mathcal{E}$] (0,2.4) to[nos, l=$K$] (2.1,2.4)
    to[R=$R$, i=$i$] (4.2,2.4) -- (5.4,2.4) to[C=$C$, v=$u$] (5.4,0) -- (0,0);
\end{circuitikz}

{\small RC 充电电路：闭合 $K$ 后，电源推动电荷
经 $R$ 流到 $C$ 的极板上。}
\end{omfigure}

\begin{theorem}[RC 电路的充电规律]\label{thm:g12:rc-rl-circuits:charging}
\index{RC 电路}
开关闭合后，电容器两端的电压满足下式，并由
$u(0) = 0$ 得出其解：
\[
RC\,\frac{du}{dt} + u = \mathcal{E},
\qquad
u(t) = \mathcal{E}\left(1 - e^{-t/\tau}\right),
\qquad \tau = RC .
\]
\end{theorem}

\begin{proof}
沿回路的电压相加（见\cref{prop:g11:circuits-and-power:laws}）：
任一时刻均有 $\mathcal{E} = u_R + u = Ri + u$；代入
$i = C\,du/dt$（见\cref{prop:g12:rc-rl-circuits:cdudt}）。现\emph{验证}
该试探解：$du/dt = (\mathcal{E}/\tau)e^{-t/\tau}$，于是
$RC\,(\mathcal{E}/RC)\,e^{-t/\tau} + \mathcal{E}(1 - e^{-t/\tau}) =
\mathcal{E}$，对所有 $t$ 成立；且 $u(0) = \mathcal{E}(1-1) = 0$。
至于除此之外没有\emph{其他}曲线满足该方程，
以及解的唯一性，将在本年度的数学课程《微分方程》中加以证明。
\end{proof}

\begin{definition}[时间常数]\label{def:g12:rc-rl-circuits:tau}
RC 电路的\emph{时间常数}\index{时间常数}为
$\tau = RC$。单位：$\unit{\ohm} \times \unit{F} =
(\unit{V/A})(\unit{C/V}) = \unit{C/A} = \unit{s}$——该电路
自身的节奏。
\end{definition}

\begin{method}[在曲线上读出 $\tau$]\label{met:g12:rc-rl-circuits:reading}
由图像直接读出，有三种等价的方法：
\begin{enumerate}
  \item \emph{原点处的切线：}$u'(0) = \mathcal{E}/\tau$，因此
        初始切线在 $t = \tau$ 处与渐近线相交；
  \item \emph{63\%：}$u(\tau) = \mathcal{E}(1 - e^{-1}) \approx
        0.63\,\mathcal{E}$；
  \item \emph{95\%：}$u(3\tau) \approx 0.95\,\mathcal{E}$——已充电，
        就实际应用而言已充满（$5\tau$ 时为 $99\%$）。
\end{enumerate}
\end{method}

\begin{proposition}[放电]\label{prop:g12:rc-rl-circuits:discharge}
电容器充电至 $U_0$ 后，仅与电阻 $R$ 相连，满足
$RC\,du/dt + u = 0$，其解为
\[
u(t) = U_0\,e^{-t/\tau}, \qquad \tau = RC
\]
——剩余 $37\%$（$\tau$ 时刻），$5\%$（$3\tau$ 时刻）。
\end{proposition}

\begin{proof}
同样由回路定律，无电源；再次代入：
$RC(-U_0/\tau)e^{-t/\tau} + U_0 e^{-t/\tau} = 0$，且 $u(0) = U_0$。
\end{proof}

\begin{omfigure}
\begin{tikzpicture}
\begin{axis}[omaxis, width=9.2cm, height=5.4cm, xlabel={$t$}, ylabel={$u$},
    xmin=0, xmax=5, ymin=0, ymax=7.4, xtick={0,1,2,3,4,5},
    xticklabels={$0$,$\tau$,$2\tau$,$3\tau$,$4\tau$,$5\tau$},
    ytick={0,2.21,3.79,6.0},
    yticklabels={$0$,$0.37\,\mathcal{E}$,$0.63\,\mathcal{E}$,$\mathcal{E}$}]
  \addplot[omExo, dashed, domain=0:5, samples=2] {6};
  \addplot[omExo, dashed, domain=0:1.18, samples=2] {6*x};
  \addplot[omExo, dashed, domain=0:1.1, samples=2] {6-6*x};
  \addplot[omDef, thick, domain=0:5, samples=80] {6*(1-exp(-x))};
  \addplot[omThm, thick, domain=0:5, samples=80] {6*exp(-x)};
  \draw[dashed, omExo] (1,0) -- (1,3.79) -- (0,3.79);
  \draw[dashed, omExo] (3,0) -- (3,5.70);
\end{axis}
\end{tikzpicture}

{\small 充电（蓝色，自 $u = 0$ 趋向 $\mathcal{E}$）与
放电（红色，自 $U_0 = \mathcal{E}$）。两条初切线均
在 $t = \tau$ 处与各自渐近线相交；$63\%$ 于 $\tau$ 时完成，$95\%$ 于 $3\tau$ 时完成。}
\end{omfigure}

\begin{example}[闪光灯的数据]\label{ex:g12:rc-rl-circuits:flashrc}
相机闪光灯的电容器 $C = \qty{800}{\micro F}$ 经
$R = \qty{1.0}{k\ohm}$: $\tau = 10^3 \times 8.0\times10^{-4} =
\qty{0.80}{s}$，故“就绪”指示灯接线设定在 $95\%$ 时动作——
需等待 $3\tau \approx \qty{2.4}{s}$：这便是等待时的嗡嗡声。
经闪光管放电时，仅有 \qty{0.50}{\ohm} 的电阻，\emph{同一}
电容器以 $\tau' = \qty{0.40}{ms}$ 放空，
比充电快两千倍。
\end{example}

\section{线圈：电感与 RL 电路}

\begin{definition}[电感器与电感]\label{def:g12:rc-rl-circuits:inductor}
称\emph{电感器}\index{电感器}为一圈导线绕成的线圈，其电流
在自身各匝中产生磁场
（见\cref{ch:g11:magnetic-fields}）；只要它\emph{变化}，线圈两端
便会出现电压：
\[
u = L\,\frac{di}{dt} .
\]
常量 $L$ 称为\emph{电感}\index{电感}，单位为
\emph{亨利}\index{亨利}（\unit{H}）：$\qty{1}{H} = \qty{1}{V.s/A}$。
\end{definition}
\admitted

\begin{remark}[线圈阻碍变化]\label{rem:g12:rc-rl-circuits:oppose}
变化的电流为何产生电压——即电磁感应——在
大学分册中推导；这里仅取定律及其性质：线圈
\emph{阻碍}其电流的变化。两者可作对照：电容器的
$u$ 不能跃变，线圈的 $i$ 也不能跃变。
\end{remark}

\begin{omfigure}
\begin{circuitikz}[font=\small]
  \draw (0,0) to[battery1, l=$\mathcal{E}$] (0,2.4) to[nos, l=$K$] (2.1,2.4)
    to[R=$R$] (4.2,2.4) -- (5.4,2.4) to[L=$L$, i=$i$, v=$u$] (5.4,0) -- (0,0);
\end{circuitikz}

{\small RL 电路：开关 $K$ 闭合后，线圈使电流
只能按自身节奏上升，时间常数 $\tau = L/R$。}
\end{omfigure}

\begin{proposition}[RL 电路中的电流上升]\label{prop:g12:rc-rl-circuits:rl}
\index{RL 电路}
闭合开关，使电源 $\mathcal{E}$、电阻器 $R$
与线圈 $L$ 串联，得 $\mathcal{E} = Ri + L\,di/dt$，即
\[
\frac{L}{R}\,\frac{di}{dt} + i = \frac{\mathcal{E}}{R},
\qquad\text{其解为}\quad
i(t) = \frac{\mathcal{E}}{R}\left(1 - e^{-t/\tau}\right),
\quad \tau = \frac{L}{R} .
\]
\end{proposition}

\begin{proof}
同样由回路定律，它与\emph{同一方程}一致，
即\cref{thm:g12:rc-rl-circuits:charging} 中的方程，令 $u \to i$，
$\mathcal{E} \to \mathcal{E}/R$，$RC \to L/R$；同一代换也验证其解，
凡是
\cref{met:g12:rc-rl-circuits:reading}中的读法均可移用。
\end{proof}

\begin{example}[电磁铁开始工作]\label{ex:g12:rc-rl-circuits:rl}
一只继电器线圈，$L = \qty{0.30}{H}$ 与 $R = \qty{6.0}{\ohm}$，接在
\qty{12}{V} 电源上：终态电流 $\mathcal{E}/R = \qty{2.0}{A}$，电流按
$\tau = 0.30/6.0 = \qty{50}{ms}$ 的节奏上升——约在 $3\tau \approx
\qty{0.15}{s}$ 时工作。
\end{example}

\begin{remark}[断开开关时的火花]\label{rem:g12:rc-rl-circuits:spark}
在通电线圈上断开开关，$i$ 必须在
不到一毫秒内骤降为零：$di/dt$ 极大，$u = L\,di/dt$ 可达数百
甚至数千伏特——火花会跨过触点。汽车点火
线圈正是有意如此，每分钟数千次。
\end{remark}

\section{储存的能量}

\begin{proposition}[电容器中的能量与线圈中的能量]\label{prop:g12:rc-rl-circuits:energy}
\index{能量!电容器中的}\index{能量!线圈中的}
电容器 $C$ 充电至电压 $u$，线圈 $L$ 通有电流
$i$ 时，所储存的能量为
\[
E_C = \tfrac12\,C u^2 = \frac{q^2}{2C},
\qquad
E_L = \tfrac12\,L i^2 .
\]
\end{proposition}

\begin{proof}
将微小电荷 $dq$ 移到电压已为 $u$ 的极板上，
所需的功为 $u\,dq$（见\cref{def:g11:circuits-and-power:voltage}）；累加这些
功，总功即直线 $u = q/C$ 之下的面积——即下图中的三角形，
因此 $E_C = \tfrac12 qu = \tfrac12 Cu^2$。对线圈，同样的三角形
（使 $i$ 增加 $di$ 需做功 $Li\,di$）给出
$E_L = \tfrac12 Li^2$。
\end{proof}

\begin{omfigure}
\begin{tikzpicture}[font=\small, x=1.35cm, y=1.05cm]
  \fill[omDef!25] (0,0) -- (4,0) -- (4,2.6) -- cycle;
  \fill[omDef!60] (2.4,0) rectangle (2.7,1.658);
  \draw[omDef, thick] (0,0) -- (4,2.6);
  \draw[-Stealth] (0,0) -- (4.8,0) node[right] {$q$};
  \draw[-Stealth] (0,0) -- (0,3.1) node[above] {$u = q/C$};
  \draw[dashed] (4,0) -- (4,2.6) -- (0,2.6);
  \node[below] at (4,0) {$q_0$}; \node[left] at (0,2.6) {$u_0$};
  \node[below, font=\footnotesize] at (2.55,0) {$dq$};
  \node[omDef!50!black, font=\footnotesize] at (2.9,0.85)
    {$E_C$};
\end{tikzpicture}

{\small 每个微元电荷 $dq$ 所需的功为 $u\,dq$；整个电荷所需的功
即三角形面积 $E_C = \tfrac12 q_0 u_0$——先移入的电荷
在低电压下便已就位。}
\end{omfigure}

\begin{example}[能量的数量级]\label{ex:g12:rc-rl-circuits:energies}
\cref{ex:g12:rc-rl-circuits:flashrc}中的闪光灯电容器：
$\tfrac12 \times 8.0\times10^{-4} \times 300^2 = \qty{36}{J}$。一个
\qty{3000}{F} 超级电容器在 \qty{2.7}{V} 下储能 \qty{1.1e4}{J}，相当于一节
五号电池的储能。\qty{10}{mH} 线圈通有 \qty{10}{A} 电流时储能 \qty{0.50}{J}。
储能少于电池，但可在毫秒内释放。
\end{example}

\begin{remark}[两类计时电路的用途]\label{rem:g12:rc-rl-circuits:applications}
\index{平滑}\index{计时电路}
“慢充快放”的特性用于驱动相机闪光灯和除颤器。
延迟 $\tau = RC$ 用于计时电路，如间歇式雨刷、
闪烁信标和楼道灯。跨接在电源两端的电容器能\emph{平滑}
电压（$u$ 不能突变）；串联线圈能平滑电流（$i$ 不能
突变）。线圈与电容器组合即构成振荡器
（见\cref{ch:g12:rlc-oscillations}）。
\end{remark}

\section{练习}

\begin{exercise}[$\star$]\label{exo:g12:rc-rl-circuits:1}
将一个 \qty{470}{\micro F} 的电容器充电至 \qty{12}{V}，求它所带的电荷；
并求使 \qty{100}{\micro F} 电容器带有 \qty{2.0}{mC} 电荷所需的电压。
\end{exercise}

\begin{exercise}[$\star$]\label{exo:g12:rc-rl-circuits:2}
将 \qty{47}{pF}、\qty{2200}{\micro F} 和 \qty{10}{F} 分别归入
电源滤波器、收音机调谐器与超级电容器中，
并计算最后一个在 \qty{2.7}{V} 时所储存的能量。
\end{exercise}

\begin{exercise}[$\star$]\label{exo:g12:rc-rl-circuits:3}
一个 \qty{220}{\micro F} 电容器两端的电压从 $0$ 均匀上升
至 \qty{5.0}{V}，历时 \qty{10}{ms}，求此时流入的电流；
电压保持不变后电流为多少，并说明理由。
\end{exercise}

\begin{exercise}[$\star$]\label{exo:g12:rc-rl-circuits:4}
求 $\tau$，当 $R = \qty{10}{k\ohm}$、$C = \qty{100}{\micro F}$ 时；
再求 $L = \qty{0.10}{H}$、$R = \qty{50}{\ohm}$ 时的值。逐项检查单位，
验证两种组合的单位都是秒。
\end{exercise}

\begin{exercise}[$\star$]\label{exo:g12:rc-rl-circuits:5}
一个电容器向 \qty{6.0}{V} 充电，电压在
$t = \qty{2.0}{s}$ 时为 \qty{3.8}{V}。读出时间常数；电荷何时达到终值的 $95\%$？
\end{exercise}

\begin{exercise}[$\star\star$]\label{exo:g12:rc-rl-circuits:6}
一个 RC 电路：$\mathcal{E} = \qty{9.0}{V}$、$R = \qty{4.7}{k\ohm}$、
$C = \qty{220}{\micro F}$。求 $\tau$、$u(\tau)$、$u(3\tau)$，以及
开关刚闭合时的电流。
\end{exercise}

\begin{exercise}[$\star\star$]\label{exo:g12:rc-rl-circuits:7}
通过代入验证 $u(t) = \mathcal{E}(1 - e^{-t/RC})$
满足 $RC\,du/dt + u = \mathcal{E}$；它在 $t = 0$ 时的值
及 $t$ 很大时的极限分别蕴含怎样的物理意义？
\end{exercise}

\begin{exercise}[$\star\star$]\label{exo:g12:rc-rl-circuits:8}
一个 \qty{100}{\micro F} 电容器充电至 \qty{12}{V} 后，经
\qty{47}{k\ohm} 的电阻放电。求 $\tau$、$u(\tau)$，以及 $u$ 降至
\qty{0.60}{V} 以下所需的时间。
\end{exercise}

\begin{exercise}[$\star\star$]\label{exo:g12:rc-rl-circuits:9}
一个 \qty{0.50}{H} 线圈中的电流从 $0$ 均匀增至
\qty{2.0}{A}，历时 \qty{40}{ms}。求线圈两端的电压；
由 $u = L\,di/dt$ 说明，为何\emph{断开}线圈电路中的开关时会产生火花，而闭合时不会？
\end{exercise}

\begin{exercise}[$\star\star$]\label{exo:g12:rc-rl-circuits:10}
\cref{ex:g12:rc-rl-circuits:rl}中的继电器：通过代入验证
$i(t) = (\mathcal{E}/R)(1 - e^{-Rt/L})$ 满足
$L\,di/dt + Ri = \mathcal{E}$，然后求 $i(\tau)$ 和 $i(3\tau)$。
\end{exercise}

\begin{exercise}[$\star\star$]\label{exo:g12:rc-rl-circuits:11}
证明 $u'(0) = \mathcal{E}/\tau$，并由此推出曲线在原点处的切线
与渐近线交于 $t = \tau$ 处。若在一条实测曲线上，
测得 $\tau = \qty{0.50}{s}$、$R = \qty{2.0}{k\ohm}$，求 $C$。
\end{exercise}

\begin{exercise}[$\star\star\star$]\label{exo:g12:rc-rl-circuits:12}
按储存能量给各器件排序，并逐一计算：一只 \qty{3000}{F} 的超级电容器，
电压为 \qty{2.7}{V}；一只闪光灯电容器，\qty{800}{\micro F}，
电压为 \qty{300}{V}；一只 \qty{10}{mH} 的线圈，电流为 \qty{10}{A}；一节 AA 电池
（$E \approx \qty{1.3e4}{J}$）。前三者能提供哪些
电池不能提供的东西？
\end{exercise}

\begin{exercise}[$\star\star\star$]\label{exo:g12:rc-rl-circuits:13}
楼梯定时灯在其充电电容器两端电压 $u < 0.63\,\mathcal{E}$ 时点亮，
在 $t = \tau$ 时熄灭。要使灯亮持续
\qty{5.0}{s}，取 $C = \qty{100}{\micro F}$，求 $R$ 应为多少？
若阈值取 $0.50\,\mathcal{E}$，证明延迟时间为 $\tau \ln 2$，并求其值。
\end{exercise}

\begin{exercise}[$\star\star\star$]\label{exo:g12:rc-rl-circuits:14}
一只 \qty{2200}{\micro F} 的电容器用于平滑电源：两次充电之间，
它单独为一个负载供电，负载电流为 \qty{0.50}{A}，持续 \qty{10}{ms}。由
$i = C\,du/dt$ 求电压降落，然后求
使降落低于 \qty{0.50}{V} 的电容。
\end{exercise}

\begin{exercise}[$\star\star\star$]\label{exo:g12:rc-rl-circuits:15}
一只点火线圈，$L = \qty{0.80}{H}$，通有 \qty{2.0}{A} 电流时，
断路器约在 \qty{1.0}{ms} 内将其切断。估算切断期间的线圈电压，
以及火花的能量；为什么缓慢切断
不会产生火花？
\end{exercise}

\section{问题：相机闪光灯}

\begin{problem}[{周末问题——相机闪光灯：两秒的
嗡鸣，一毫秒的辉煌，壁挂式的“近亲”，能重启
心脏——一个指数函数统领全部}]
\label{pb:g12:rc-rl-circuits:1}
一台小型闪光灯将能量储存在一只电容器中，
$C = \qty{800}{\micro F}$，经 $R = \qty{1.0}{k\ohm}$ 充电，电源为
视为理想电源的转换器，$\mathcal{E} = \qty{300}{V}$；
“就绪”指示灯在满电压的 $95\%$ 时亮起。闪光时，
电容器仅与闪光管相连，$R_f = \qty{0.50}{\ohm}$。已知：一节
AA 电池电压为 \qty{1.5}{V}，容量为 \qty{2.4}{A.h}；$e = \qty{1.60e-19}{C}$。

\medskip
\noindent\textbf{第一部分——能量库。}
\begin{enumerate}
  \item 计算电压为 \qty{300}{V} 时电容器上的电荷量。
  \item 这相当于多少个元电荷？
  \item 计算储存的能量 $E_C$。
  \item AA 电池储存的电荷量为 $q = It \approx \qty{8.6e3}{C}$，因此约
        \qty{1.3e4}{J}。这是电容器能量的多少倍？
  \item 那么何必使用电容器？用一句话作答，并使用“功率”一词。
\end{enumerate}

\medskip
\noindent\textbf{第二部分——等待绿灯。}
\begin{enumerate}[resume]
  \item 写出回路定律，并将其化为
        关于 $u(t)$ 的微分方程。
  \item 代入验证 $u(t) = \mathcal{E}(1 - e^{-t/RC})$
        是否满足该方程，初始条件为 $u(0) = 0$。
  \item 计算时间常数 $\tau$。
  \item 经过多长时间“就绪”指示灯变为绿色，此时
        电压为多少？并说明两秒嗡鸣的来源。
  \item 计算初始充电电流；为什么它之后会
        逐渐衰减？
\end{enumerate}

\medskip
\noindent\textbf{第三部分——一毫秒的辉煌。}
\begin{enumerate}[resume]
  \item 写出经闪光管放电的规律 $u(t)$，并计算
        新的时间常数 $\tau'$。
  \item 计算通过闪光管的峰值电流。
  \item 计算峰值功率，并将其与一只 \qty{2.2}{kW} 的电水壶作比较。
  \item 取闪光在 $3\tau'$ 时结束，求其持续时间与
        平均功率。
  \item 求充电时间与闪光时间之比：同一电容器，
        同样的电荷量——有何变化，这使电容器成为
        一个能量储存器，还是一根功率杠杆？
\end{enumerate}

\medskip
\noindent\textbf{第四部分——壁挂式的“近亲”。}
除颤器须释放 $E_C = \qty{200}{J}$ 的能量，所用电压为
$U_0 = \qty{1500}{V}$；胸部与电极板的电阻约为 \qty{50}{\ohm}。
\begin{enumerate}[resume]
  \item 计算所需的电容。
  \item 计算储存的电荷量。
  \item 计算放电的时间常数与持续时间（$3\tau$）；
        并与闪光灯作比较。
  \item 计算患者体内的峰值电流与峰值功率。
  \item 再充电电路须使设备在
        \qty{6.0}{s} 内达到就绪（$95\%$）。求所需的充电电阻和
        初始充电电流——设计者据此订购器件所需的数值。
\end{enumerate}
\end{problem}
